\par\Needspace{6\baselineskip}\medskip
\begin{center}\hypertarget{native-183-article137}{}\pdfbookmark[2]{137.}{article-137}137.\end{center}
Negotium \emph{secundum} erit determinatio situs relativi illorum trium circulorum
maximorum inter se, qui pendebit a situ intersectionum mutuarum et ab
inclinationibus. Quae si absque ambiguitate ad notiones claras ac generales
reducere cupimus, ita ut non opus sit pro singulis casibus diversis ad figuras
peculiares recurrere, quasdam dilucidationes praeliminares praemittere oportebit.
\emph{Primo} scilicet in quovis circulo maximo duae \emph{directiones} oppositae aliquo modo
distinguendae sunt, quod fiet, dum alteram tamquam progressivam seu positivam,
alteram tamquam retrogradam seu negativam consideramus. Quod quum
per se prorsus arbitrarium sit, ut normam certam stabiliamus, semper directiones
ab $A,A',A''$ versus $B,B',B''$ ceu positivas considerabimus; ita e.\,g. si
intersectio circuli primi cum secundo per distantiam positivam a puncto $A$
exhibetur, haec capienda subintelligetur ab $A$ versus $B$ (ut $D''$ in figura nostra);
si vero negativa esset, ipsam ab altera parte ipsius $A$ sumere oporteret. \emph{Secundo}
vero etiam duo haemisphaeria, in quae omnis circulus maximus sphaeram
integram dirimit, denominationibus idoneis distinguenda sunt: et quidem
hemisphaerium \emph{superius} vocabimus, quod in superficie interiori sphaerae circulum
maximum directione progressiva permeanti ad dextram est, alterum \emph{inferius}.
Plaga itaque superior analoga erit hemisphaerio boreali respectu eclipticae
vel aequatoris, inferior australi.

His rite intellectis, \emph{ambas} duorum circulorum maximorum intersectiones
commode ab invicem distinguere licebit: in una scilicet circulus primus e secundi
regione inferiori in superiorem tendit, vel quod idem est secundus e
primi regione superiori in inferiorem; in altera intersectione opposita locum
habent. Per se quidem prorsus arbitrarium est, quasnam intersectiones in
problemate nostro eligere velimus: sed ut hic quoque iuxta normam invariabilem
procedamus, eas semper adoptabimus ($D,D',D''$ in fig. 4), ubi resp.
circulus tertius $A''B''$ in secundi $A'B'$, tertius in primi $AB$, secundus in primi
plagam superiorem transit. Situs harum intersectionum determinabitur per
\hypertarget{native-184}{}ipsarum distantias a punctis $A'$ et $A''$, $A$ et $A''$, $A$ et $A'$, quas simpliciter
per $A'D,A''D,AD',A''D',AD'',A'D''$ designabimus. Quibus ita factis circulorum
inclinationes mutuae erunt anguli, qui resp. in his intersectionum
punctis $D,D',D''$ inter circulorum se secantium partes eas continentur, quae
in directione progressiva iacent: has inclinationes, semper inter $0$ et $180^0$ accipiendas,
per $\epsilon,\epsilon',\epsilon''$ denotabimus. Determinatio harum novem quantitatum
incognitarum e cognitis manifesto ab eodem problemate pendet, quod in art. 55
tractavimus: habemus itaque aequationes sequentes:
\[
\begin{aligned}
3)\quad&\sin\tfrac12\epsilon\sin\tfrac12(A'D+A''D)=\sin\tfrac12(l''-l')\sin\tfrac12(\gamma''+\gamma')\\[.5ex]
4)\quad&\sin\tfrac12\epsilon\cos\tfrac12(A'D+A''D)=\cos\tfrac12(l''-l')\sin\tfrac12(\gamma''-\gamma')\\[.5ex]
5)\quad&\cos\tfrac12\epsilon\sin\tfrac12(A'D-A''D)=\sin\tfrac12(l''-l')\cos\tfrac12(\gamma''+\gamma')\\[.5ex]
6)\quad&\cos\tfrac12\epsilon\cos\tfrac12(A'D-A''D)=\cos\tfrac12(l''-l')\cos\tfrac12(\gamma''-\gamma').
\end{aligned}
\]
Ex aequationibus 3 et 4 innotescent $\tfrac12(A'D+A''D)$ et $\sin\tfrac12\epsilon$, e duabus reliquis
$\tfrac12(A'D-A''D)$ et $\cos\tfrac12\epsilon$; hinc $A'D,A''D$ et $\epsilon$. Ambiguitas determinationi arcuum
$\tfrac12(A'D+A''D)$, $\tfrac12(A'D-A''D)$ per tangentes adhaerens conditione ea
decidetur, quod $\sin\tfrac12\epsilon$ et $\cos\tfrac12\epsilon$ positivi evadere debent, consensusque inter
$\sin\tfrac12\epsilon$ et $\cos\tfrac12\epsilon$ toti calculo confirmando inserviet.

Determinatio quantitatum $AD',A''D',\epsilon',AD'',A'D'',\epsilon''$ prorsus simili
modo perficietur, neque opus erit octo aequationes ad hunc calculum adhibendas
huc transscribere, quippe quae ex aequ. 3---6 sponte prodeunt, si
\[
\begin{array}{r c|c|c|c|c|c}
 & A'D & A''D & \epsilon & l''-l' & \gamma'' & \gamma'\\
\text{cum} & AD' & A''D' & \epsilon' & l''-l & \gamma'' & \gamma\\
\text{vel cum} & AD'' & A'D'' & \epsilon'' & l'-l & \gamma' & \gamma
\end{array}
\]
resp. commutantur.

Nova adhuc totius calculi confirmatio derivari potest e relatione mutua
inter latera angulosque trianguli sphaerici inter puncta $D,D',D''$ formati,
unde demanant aequationes generalissime verae, quamcunque situm haec puncta
habeant:
\[
\frac{\sin(AD'-AD'')}{\sin\epsilon}=\frac{\sin(A'D-A'D'')}{\sin\epsilon'}=\frac{\sin(A''D-A''D')}{\sin\epsilon''}\mathpunct{\cdot}
\]

Denique si loco eclipticae aequator tamquam planum fundamentale electus
\hypertarget{native-185}{}est, calculus mutationem non subit, nisi quod pro terrae locis heliocentricis
$A,A',A''$ substituere oportet ea aequatoris puncta, ubi a circulis $AB,A'B',A''B''$
secatur; accipiendae sunt itaque pro $l,l',l''$ ascensiones rectae harum
intersectionum, nec non pro $A'D$ distantia puncti $D$ ab intersectione secunda etc.
